% ========================================================
% BÀI TẬP TOÁN 5 - KẾT NỐI TRI THỨC VỚI CUỘC SỐNG
% Bài 52: Thể tích của hình hộp chữ nhật
% Định dạng: BTTeX Part III - Trắc nghiệm điền số ngắn
% ========================================================

% Câu 1
\begin{ex}
  Một hình hộp chữ nhật có chiều dài $6\text{ cm}$, chiều rộng $4\text{ cm}$ và chiều cao $3\text{ cm}$. Thể tích của hình hộp chữ nhật đó là bao nhiêu xăng-ti-mét khối?
  \shortans{72}
  \loigiai{
    Thể tích là: 6 x 4 x 3 = 72 cm3.
  }
\end{ex}

% Câu 2
\begin{ex}
  Một hình hộp chữ nhật có chiều dài $7\text{ cm}$, chiều rộng $4\text{ cm}$ và chiều cao $3\text{ cm}$. Thể tích của hình hộp chữ nhật đó là bao nhiêu xăng-ti-mét khối?
  \shortans{84}
  \loigiai{
    Thể tích là: 7 x 4 x 3 = 84 cm3.
  }
\end{ex}

% Câu 3
\begin{ex}
  Một hình hộp chữ nhật có chiều dài $8\text{ cm}$, chiều rộng $4\text{ cm}$ và chiều cao $3\text{ cm}$. Thể tích của hình hộp chữ nhật đó là bao nhiêu xăng-ti-mét khối?
  \shortans{96}
  \loigiai{
    Thể tích là: 8 x 4 x 3 = 96 cm3.
  }
\end{ex}

% Câu 4
\begin{ex}
  Một hình hộp chữ nhật có chiều dài $9\text{ cm}$, chiều rộng $4\text{ cm}$ và chiều cao $3\text{ cm}$. Thể tích của hình hộp chữ nhật đó là bao nhiêu xăng-ti-mét khối?
  \shortans{108}
  \loigiai{
    Thể tích là: 9 x 4 x 3 = 108 cm3.
  }
\end{ex}

% Câu 5
\begin{ex}
  Một hình hộp chữ nhật có chiều dài $10\text{ cm}$, chiều rộng $4\text{ cm}$ và chiều cao $3\text{ cm}$. Thể tích của hình hộp chữ nhật đó là bao nhiêu xăng-ti-mét khối?
  \shortans{120}
  \loigiai{
    Thể tích là: 10 x 4 x 3 = 120 cm3.
  }
\end{ex}

% Câu 6
\begin{ex}
  Một hình hộp chữ nhật có chiều dài $11\text{ cm}$, chiều rộng $4\text{ cm}$ và chiều cao $3\text{ cm}$. Thể tích của hình hộp chữ nhật đó là bao nhiêu xăng-ti-mét khối?
  \shortans{132}
  \loigiai{
    Thể tích là: 11 x 4 x 3 = 132 cm3.
  }
\end{ex}

% Câu 7
\begin{ex}
  Một hình hộp chữ nhật có chiều dài $12\text{ cm}$, chiều rộng $4\text{ cm}$ và chiều cao $3\text{ cm}$. Thể tích của hình hộp chữ nhật đó là bao nhiêu xăng-ti-mét khối?
  \shortans{144}
  \loigiai{
    Thể tích là: 12 x 4 x 3 = 144 cm3.
  }
\end{ex}

% Câu 8
\begin{ex}
  Một hình hộp chữ nhật có chiều dài $13\text{ cm}$, chiều rộng $4\text{ cm}$ và chiều cao $3\text{ cm}$. Thể tích của hình hộp chữ nhật đó là bao nhiêu xăng-ti-mét khối?
  \shortans{156}
  \loigiai{
    Thể tích là: 13 x 4 x 3 = 156 cm3.
  }
\end{ex}

% Câu 9
\begin{ex}
  Một hình hộp chữ nhật có chiều dài $14\text{ cm}$, chiều rộng $4\text{ cm}$ và chiều cao $3\text{ cm}$. Thể tích của hình hộp chữ nhật đó là bao nhiêu xăng-ti-mét khối?
  \shortans{168}
  \loigiai{
    Thể tích là: 14 x 4 x 3 = 168 cm3.
  }
\end{ex}

% Câu 10
\begin{ex}
  Một hình hộp chữ nhật có chiều dài $15\text{ cm}$, chiều rộng $4\text{ cm}$ và chiều cao $3\text{ cm}$. Thể tích của hình hộp chữ nhật đó là bao nhiêu xăng-ti-mét khối?
  \shortans{180}
  \loigiai{
    Thể tích là: 15 x 4 x 3 = 180 cm3.
  }
\end{ex}

% Câu 11
\begin{ex}
  Một hình hộp chữ nhật có chiều dài $16\text{ cm}$, chiều rộng $4\text{ cm}$ và chiều cao $3\text{ cm}$. Thể tích của hình hộp chữ nhật đó là bao nhiêu xăng-ti-mét khối?
  \shortans{192}
  \loigiai{
    Thể tích là: 16 x 4 x 3 = 192 cm3.
  }
\end{ex}

% Câu 12
\begin{ex}
  Một hình hộp chữ nhật có chiều dài $17\text{ cm}$, chiều rộng $4\text{ cm}$ và chiều cao $3\text{ cm}$. Thể tích của hình hộp chữ nhật đó là bao nhiêu xăng-ti-mét khối?
  \shortans{204}
  \loigiai{
    Thể tích là: 17 x 4 x 3 = 204 cm3.
  }
\end{ex}

% Câu 13
\begin{ex}
  Một hình hộp chữ nhật có chiều dài $18\text{ cm}$, chiều rộng $4\text{ cm}$ và chiều cao $3\text{ cm}$. Thể tích của hình hộp chữ nhật đó là bao nhiêu xăng-ti-mét khối?
  \shortans{216}
  \loigiai{
    Thể tích là: 18 x 4 x 3 = 216 cm3.
  }
\end{ex}

% Câu 14
\begin{ex}
  Một hình hộp chữ nhật có chiều dài $19\text{ cm}$, chiều rộng $4\text{ cm}$ và chiều cao $3\text{ cm}$. Thể tích của hình hộp chữ nhật đó là bao nhiêu xăng-ti-mét khối?
  \shortans{228}
  \loigiai{
    Thể tích là: 19 x 4 x 3 = 228 cm3.
  }
\end{ex}

% Câu 15
\begin{ex}
  Một hình hộp chữ nhật có chiều dài $20\text{ cm}$, chiều rộng $4\text{ cm}$ và chiều cao $3\text{ cm}$. Thể tích của hình hộp chữ nhật đó là bao nhiêu xăng-ti-mét khối?
  \shortans{240}
  \loigiai{
    Thể tích là: 20 x 4 x 3 = 240 cm3.
  }
\end{ex}

% Câu 16
\begin{ex}
  Một hình hộp chữ nhật có chiều dài $21\text{ cm}$, chiều rộng $4\text{ cm}$ và chiều cao $3\text{ cm}$. Thể tích của hình hộp chữ nhật đó là bao nhiêu xăng-ti-mét khối?
  \shortans{252}
  \loigiai{
    Thể tích là: 21 x 4 x 3 = 252 cm3.
  }
\end{ex}

% Câu 17
\begin{ex}
  Một hình hộp chữ nhật có chiều dài $22\text{ cm}$, chiều rộng $4\text{ cm}$ và chiều cao $3\text{ cm}$. Thể tích của hình hộp chữ nhật đó là bao nhiêu xăng-ti-mét khối?
  \shortans{264}
  \loigiai{
    Thể tích là: 22 x 4 x 3 = 264 cm3.
  }
\end{ex}

% Câu 18
\begin{ex}
  Một hình hộp chữ nhật có chiều dài $23\text{ cm}$, chiều rộng $4\text{ cm}$ và chiều cao $3\text{ cm}$. Thể tích của hình hộp chữ nhật đó là bao nhiêu xăng-ti-mét khối?
  \shortans{276}
  \loigiai{
    Thể tích là: 23 x 4 x 3 = 276 cm3.
  }
\end{ex}

% Câu 19
\begin{ex}
  Một hình hộp chữ nhật có chiều dài $24\text{ cm}$, chiều rộng $4\text{ cm}$ và chiều cao $3\text{ cm}$. Thể tích của hình hộp chữ nhật đó là bao nhiêu xăng-ti-mét khối?
  \shortans{288}
  \loigiai{
    Thể tích là: 24 x 4 x 3 = 288 cm3.
  }
\end{ex}

% Câu 20
\begin{ex}
  Một hình hộp chữ nhật có chiều dài $25\text{ cm}$, chiều rộng $4\text{ cm}$ và chiều cao $3\text{ cm}$. Thể tích của hình hộp chữ nhật đó là bao nhiêu xăng-ti-mét khối?
  \shortans{300}
  \loigiai{
    Thể tích là: 25 x 4 x 3 = 300 cm3.
  }
\end{ex}

% Câu 21
\begin{ex}
  Một hình hộp chữ nhật có chiều dài $26\text{ cm}$, chiều rộng $4\text{ cm}$ và chiều cao $3\text{ cm}$. Thể tích của hình hộp chữ nhật đó là bao nhiêu xăng-ti-mét khối?
  \shortans{312}
  \loigiai{
    Thể tích là: 26 x 4 x 3 = 312 cm3.
  }
\end{ex}

% Câu 22
\begin{ex}
  Một hình hộp chữ nhật có chiều dài $27\text{ cm}$, chiều rộng $4\text{ cm}$ và chiều cao $3\text{ cm}$. Thể tích của hình hộp chữ nhật đó là bao nhiêu xăng-ti-mét khối?
  \shortans{324}
  \loigiai{
    Thể tích là: 27 x 4 x 3 = 324 cm3.
  }
\end{ex}

% Câu 23
\begin{ex}
  Một hình hộp chữ nhật có chiều dài $28\text{ cm}$, chiều rộng $4\text{ cm}$ và chiều cao $3\text{ cm}$. Thể tích của hình hộp chữ nhật đó là bao nhiêu xăng-ti-mét khối?
  \shortans{336}
  \loigiai{
    Thể tích là: 28 x 4 x 3 = 336 cm3.
  }
\end{ex}

% Câu 24
\begin{ex}
  Một hình hộp chữ nhật có chiều dài $29\text{ cm}$, chiều rộng $4\text{ cm}$ và chiều cao $3\text{ cm}$. Thể tích của hình hộp chữ nhật đó là bao nhiêu xăng-ti-mét khối?
  \shortans{348}
  \loigiai{
    Thể tích là: 29 x 4 x 3 = 348 cm3.
  }
\end{ex}

% Câu 25
\begin{ex}
  Một hình hộp chữ nhật có chiều dài $30\text{ cm}$, chiều rộng $4\text{ cm}$ và chiều cao $3\text{ cm}$. Thể tích của hình hộp chữ nhật đó là bao nhiêu xăng-ti-mét khối?
  \shortans{360}
  \loigiai{
    Thể tích là: 30 x 4 x 3 = 360 cm3.
  }
\end{ex}
